% Source: https://artofproblemsolving.com/community/c3h1732120_simplification_problem
% Stable identity: aops-topic-1732120-q01
First, $1+a^2+a^4=(a^2+1)^2-a^2=(a^2+a+1)(a^2-a+1)$, and both factors are positive. Look for $u,v\ge0$ with $(\sqrt u+\sqrt v)^2=u+v+2\sqrt{uv}$ equal to the expression under the outer root, i.e.
\[
u+v=1+a^2,\qquad 4uv=1+a^2+a^4 .
\]
Since $(a^2+a+1)+(a^2-a+1)=2(1+a^2)$, the choice $u=\dfrac{a^2+a+1}{2}$, $v=\dfrac{a^2-a+1}{2}$ works: $u+v=1+a^2$ and $4uv=(a^2+a+1)(a^2-a+1)=1+a^2+a^4$. Therefore
\[
\sqrt{1+a^2+\sqrt{1+a^2+a^4}}=\sqrt{\dfrac{a^2+a+1}{2}}+\sqrt{\dfrac{a^2-a+1}{2}}=\boxed{\dfrac{\sqrt{2a^2+2a+2}+\sqrt{2a^2-2a+2}}{2}}.
\]
(Check at $a=0$: both sides equal $\sqrt2$; at $a=1$: $\sqrt{2+\sqrt3}=\dfrac{\sqrt6+\sqrt2}{2}$.)
